\chapter{Risque opérationnel et fraude de trader}\label{ch:rc:operational-risk-and-rogue-trading}

Le 26 février 1995, la banque londonienne Barings s'effondre. Son trader à Singapour avait accumulé des pertes sur un compte numéroté 88888 que Londres n'a jamais vu~: plus de
20 millions de livres à la fin de 1993, plus de 200 millions à la fin de 1994, et 827 millions à
la faillite. Il était chargé à la fois du trading et du back office qui l'enregistrait~; les auditeurs internes de la banque
avaient recommandé en 1994 qu'il ne le soit pas~; Londres lui avait envoyé plus de 300
millions de livres à prêter à des clients qui avaient déposé environ 31 millions de collatéral, et personne ne demanda
pourquoi. Aucun des modèles des chapitres précédents n'aurait mesuré ce risque~: les positions étaient
cachées à ces modèles. Ce chapitre porte sur le risque qui échappe aux modèles~: le \omterm{def:rc:operational-risk-and-rogue-trading:oprisk}{risque opérationnel}, et
sa forme la plus coûteuse dans le trading, les positions non autorisées dissimulées par ceux qui les gèrent.

\section{Risque opérationnel et sa mesure}

\begin{definition}[Risque opérationnel, événement de perte opérationnelle]\label{def:rc:operational-risk-and-rogue-trading:oprisk}
Le \emph{risque opérationnel}\index{risque opérationnel} est le risque de perte résultant de processus internes
inadéquats ou défaillants, de personnes et de systèmes, ou d'événements externes~; il inclut le risque juridique et
exclut le risque stratégique et le risque de réputation. Un \emph{événement de perte opérationnelle}\index{evenement@événement de perte opérationnelle}
est la survenue d'une telle perte, consignée avec sa date, son montant, sa ligne de métier et sa cause.
\end{definition}

Les pertes opérationnelles sont rares et à queue épaisse~: la plupart sont de petites erreurs, quelques-unes sont catastrophiques. Les banques
les recensent dans des bases de pertes, les rattachent à leurs causes et aux contrôles, et détiennent des fonds propres pour les couvrir.

\begin{method}[L'approche standard de Bâle]\label{met:rc:operational-risk-and-rogue-trading:sa}
Calculer l'indicateur d'activité (BI), une mesure de la taille de la banque à partir de son compte de résultat. La
composante de l'indicateur d'activité (BIC) applique des coefficients marginaux de 12~\% jusqu'à 1~milliard d'euros de BI, 15~\%
de 1 à 30 milliards d'euros et 18~\% au-delà. Il est ensuite multiplié par le multiplicateur de perte interne
$\mathrm{ILM} = \ln\bigl(e-1+(\mathrm{LC}/\mathrm{BIC})^{0.8}\bigr)$, où la composante de perte LC vaut 15 fois la
perte opérationnelle annuelle moyenne des dix dernières années. Le capital vaut $\mathrm{BIC}\times\mathrm{ILM}$.
\end{method}

\begin{example}[Capital au titre du risque opérationnel]\label{ex:rc:operational-risk-and-rogue-trading:capital}
Une banque dont le BI vaut 35~milliards d'euros a un BIC de $0.12+29\times0.15+5\times0.18 = 5.37$ milliards d'euros. Avec des pertes
annuelles moyennes de 0.5 milliard d'euros, sa composante de perte vaut 7.5 milliards d'euros, son ILM 1.107 et son capital 5.94 milliards d'euros~;
avec 0.25 milliard de pertes, l'ILM vaut 0.904 et le capital 4.85 milliards d'euros.
\end{example}

\section{Les cas documentés}

\begin{definition}[Fraude de trader, opération fictive]\label{def:rc:operational-risk-and-rogue-trading:rogue}
La \emph{fraude de trader}\index{fraude de trader} (rogue trading) est une activité de trading menée au-delà des pouvoirs d'un trader, en volume ou en nature,
dissimulée à l'entreprise. Une \emph{opération fictive}\index{opération fictive} est une écriture sans aucune réalité économique,
passée pour compenser ou maquiller une position réelle dans les registres de l'entreprise.
\end{definition}

\begin{dated}{2026-09}{Six cas}\label{dat:rc:operational-risk-and-rogue-trading:cases}
Barings (1995)~: pertes de 827~millions de livres sterling dissimulées sur le compte 88888 par un trader qui dirigeait aussi le
back office (déclaration du chancelier sur le rapport du Board of Banking Supervision). Daiwa (1995)~:
un trader new-yorkais a perdu environ 1.1~milliard de dollars sur des titres du Trésor en onze ans, en couvrant les
pertes en vendant des titres détenus en conservation et en falsifiant les registres~; il dirigeait aussi le service de conservation
(ordonnance de la Federal Reserve). Sumitomo (1995--96)~: la CFTC américaine a constaté qu'après de lourdes pertes liées au
trading spéculatif, le principal trader sur le cuivre de la société s'est livré à une manœuvre visant à manipuler les
cours du cuivre~; la société a versé 150~millions de dollars et, en septembre 1996, a chiffré ses pertes à 2.6~milliards de dollars.
Allfirst (1997--2002)~: pertes de 691.2~millions de dollars dissimulées au moyen d'options de change fictives qui
semblaient couvrir des positions réelles (rapport annuel de la maison mère). Société Générale (2008)~: une position
longue de 49~milliards d'euros sur des contrats à terme sur indices, constituée du 2 au 18 janvier, a coûté 6.4~milliards d'euros à dénouer,
soit 4.9~milliards d'euros après déduction des gains de 2007~; elle a été masquée par 947 \omterm{def:rc:operational-risk-and-rogue-trading:rogue}{opérations fictives} annulées avant confirmation
(rapport d'inspection de la banque). UBS (2011)~: pertes de 2.3~milliards de dollars sur un desk ETF, dissimulées par des
comptabilisations tardives, transactions internes non appariées et fausses transactions à règlement différé~; le régulateur britannique a infligé
à la banque une amende de 29.7~millions de livres sterling.
\end{dated}

\section{Les défaillances de contrôle communes}

\begin{definition}[Séparation des tâches, congés obligatoires consécutifs]\label{def:rc:operational-risk-and-rogue-trading:sod}
La \emph{séparation des tâches}\index{séparation des tâches} (segregation of duties) sépare les personnes qui négocient de celles qui
enregistrent, confirment, règlent, valorisent et déclarent les transactions. Les \emph{congés obligatoires consécutifs}\index{congés obligatoires consécutifs} (block leave) obligent chaque
trader à s'éloigner du desk, sans accès à ses systèmes, pendant une période consécutive chaque année, de sorte que
les positions qui exigent une attention quotidienne pour rester cachées soient exposées.
\end{definition}

Les défaillances documentées présentent un même schéma (la \cref{fig:rc:operational-risk-and-rogue-trading:chain})~: un
trader qui contrôle l'enregistrement ou la conservation de ses propres transactions (Barings, Daiwa)~; des contreparties internes
et des transactions que personne n'apparie ni ne confirme (UBS), ou des confirmations que le trader falsifie (Allfirst)~; des \omterm{def:rc:operational-risk-and-rogue-trading:rogue}{opérations fictives}
à dates de valeur lointaines, annulées avant d'être confirmées (Société Générale)~; un financement
que personne ne rapproche des positions qu'il finance (Barings). Chacun est un maillon de la chaîne de contrôles
que la transaction devait traverser, et chacun est rompu par une seule personne.

\begin{omfigure}
\begin{tikzpicture}[font=\small,box/.style={draw,rounded corners,minimum width=2.3cm,minimum height=0.95cm,align=center,text width=2.2cm},
  note/.style={omThm,align=center,font=\scriptsize,text width=2.6cm,anchor=north}]
\node[box,fill=omThm!10] (t) at (0,0) {transaction};
\node[box] (b) at (3,0) {saisie};
\node[box] (c) at (6,0) {confirmation};
\node[box] (s) at (9,0) {règlement-livraison et espèces};
\node[box] (p) at (12,0) {P\&L explain et IPV};
\draw[->,thick] (t) -- (b);
\draw[->,thick] (b) -- (c);
\draw[->,thick] (c) -- (s);
\draw[->,thick] (s) -- (p);
\node[note] at (3,-0.75) {saisies tardives~; absence de filtre hors marché (UBS)};
\node[note] at (6,-0.75) {transactions internes non appariées (UBS)};
\node[note] at (9,-0.75) {règlement différé (UBS)~; financement non rapproché (Barings)};
\node[note] at (12,-0.75) {pertes dans un compte caché (Barings)};
\node[omDef,align=center,font=\footnotesize] at (6,1.0) {une seule personne contrôlant le trading et le back office rompt tous les maillons (Barings)};
\end{tikzpicture}

\omcaption{La chaîne de contrôles que traverse une transaction, et l'endroit où les cas documentés l'ont rompue.
Chaque maillon constitue une preuve indépendante que la position existe et qu'elle est bien ce qu'elle prétend être. Schéma de principe~;
cas tirés de la déclaration du chancelier sur Barings (1995) et de l'avis de la FSA de 2012 adressé à UBS.}\label{fig:rc:operational-risk-and-rogue-trading:chain}
\end{omfigure}

\section{Détecter les opérations fictives et dissimulées}

La surveillance traduit ce schéma en règles portant sur les registres de transactions~: une transaction enregistrée longtemps après
son exécution~; annulations et modifications autour des dates de déclaration, lorsque les positions sont mesurées~; prix très
éloignés du marché~; transactions internes sans transaction correspondante dans l'autre portefeuille~; dates de règlement-livraison
très lointaines dans le futur, pour des transactions jamais destinées à être réglées.

\omcode{code/firm/tradecontrol/firm_tradecontrol.py}{36}{55}{Cinq règles de surveillance sur un blotter de transactions.}\label{lst:rc:operational-risk-and-rogue-trading:rules}

\begin{example}[Règles sur un blotter synthétique]\label{ex:rc:operational-risk-and-rogue-trading:rules}
Un blotter synthétique contient 2\,000 transactions authentiques, avec le bruit des opérations réelles (4~\% saisies en retard, quelques
exécutions hors marché, des transactions internes dont 2~\% n'ont pas de contrepartie miroir, des annulations à des moments aléatoires), et 60 transactions
d'un système de dissimulation. Isolément, chaque règle détecte 23~\% à 52~\% des transactions du système, avec des taux de fausses alertes
allant de 0.25~\% (transactions internes sans correspondance) à 4.4~\% (saisies tardives). Signaler toute transaction qui enfreint une règle
détecte 88.3~\% du système pour un taux de fausses alertes de 7.1~\%, soit 142 transactions authentiques à examiner~; exiger deux
règles détecte 61.7~\% du système avec une seule fausse alerte (\cref{fig:rc:operational-risk-and-rogue-trading:rules}).
\end{example}

\begin{omfigure}
\begin{tikzpicture}
\begin{axis}[width=11.5cm,height=5.6cm,ybar,bar width=8pt,
  symbolic x coords={late booking,cancel near report,off-market,internal unmatched,deferred settle,any rule,two rules},xtick=data,xticklabels={{saisie tardive},{annulation avant reporting},{hors marché},{interne non rapprochée},{règlement différé},{une règle},{deux règles}},
  x tick label style={font=\footnotesize,rotate=35,anchor=north east},
  ylabel={\% des transactions signalées},
  tick label style={font=\small},label style={font=\small},
  ymin=0,ymax=100,grid=major,grid style={gray!25},
  legend columns=2,legend style={font=\footnotesize,at={(0.5,1.03)},anchor=south,draw=none,fill=none,/tikz/every even column/.append style={column sep=8pt}},
  table/col sep=comma]
\addplot[fill=omDef!60,draw=omDef] table[x=rule,y=hit]{figdata/rates-credit-risk/28-operational-risk-and-rogue-trading/rules.csv};
\addplot[fill=omThm!60,draw=omThm] table[x=rule,y=false]{figdata/rates-credit-risk/28-operational-risk-and-rogue-trading/rules.csv};
\legend{transactions du système détectées (taux de détection),transactions authentiques signalées (fausses alertes)}
\end{axis}
\end{tikzpicture}

\omcaption{Taux de détection et taux de fausses alertes des règles de surveillance sur le blotter synthétique, isolément et
combinées. Aucune règle isolée ne détecte la plupart des transactions du système~; réunies, elles en détectent la plus grande part, au prix d'alertes à instruire.
Données~: synthétiques~; tutoriel du chapitre.}\label{fig:rc:operational-risk-and-rogue-trading:rules}
\end{omfigure}

Une alerte ne vaut que par son instruction~: une alerte classée sur la seule explication de la personne
qu'elle concerne ne contrôle rien. Un journal d'audit en ajout seul, dans lequel chaque entrée porte le hachage de la précédente,
rend le registre des transactions et de leur instruction infalsifiable~: une entrée modifiée rompt la chaîne.

\section{Tutoriel~: surveillance sur un blotter}\label{tut:rc:operational-risk-and-rogue-trading:tutorial}

\begin{tutorial}
\textbf{Objectif.} Exécuter les cinq règles sur un blotter synthétique comportant un système dissimulé, mesurer les taux de détection et de fausses alertes,
tenir un journal infalsifiable et calculer le capital de \omterm{def:rc:operational-risk-and-rogue-trading:oprisk}{risque opérationnel}. \textbf{État final~:} les résultats figurant dans les
\cref{ex:rc:operational-risk-and-rogue-trading:capital,ex:rc:operational-risk-and-rogue-trading:rules}
et le graphique.
\begin{tutsteps}
\item \textbf{Blotter}~: \texttt{blotter()}, transactions authentiques et du système, avec leurs étiquettes.
\item \textbf{Règles}~: \texttt{run\_rules}, \texttt{evaluate} avec une règle puis deux règles exigées.
\item \textbf{Audit}~: \texttt{audit\_demo()}~: la chaîne se vérifie, puis échoue après une modification.
\item \textbf{Capital}~: \texttt{op\_capital(35, 0.5)}~; \texttt{fig\_rc\_tradecontrol.py}.
\end{tutsteps}
\textbf{Pour aller plus loin.} Pondérer les règles par leurs taux de fausses alertes~; ajouter une règle sur les profits jugés trop élevés
au regard du risque du desk (un ratio de Sharpe supérieur à un seuil).
\end{tutorial}

\section{Construction~: contrôle des transactions}\label{bld:rc:operational-risk-and-rogue-trading:tradecontrol}

\begin{build}
\bfield{Rôle} La surveillance par la banque de ses propres transactions~: règles appliquées à chaque saisie, notation des alertes et
un journal d'audit infalsifiable pour les transactions, les modifications et les instructions.
\bfield{Interface} \texttt{Trade}~; \texttt{late\_booking}, \texttt{cancel\_amend\_near}, \texttt{off\_market},
\texttt{internal\_unmatched}, \texttt{deferred\_settlement}~; \texttt{run\_rules}, \texttt{evaluate}~;
\texttt{AuditLog} avec \texttt{append} et \texttt{verify}~; \texttt{trade\_event}.
\bfield{Règles} Les règles sont indépendantes et peu coûteuses~; les scores s'additionnent~; les alertes ne se clôturent qu'avec des éléments probants consignés
dans le journal.
\bfield{Tests d'acceptation} \texttt{code/firm/tradecontrol/tests/}~: chaque règle signale son cas planté et non
le cas propre~; l'arithmétique de l'évaluation~; la chaîne d'audit se vérifie et détecte une modification.
\bfield{Pour aller plus loin} Règles sur le P\&L et la trésorerie (profits au regard du risque, marge au regard des positions)~; références
de comportement par trader~; liens avec les confirmations et les données de règlement-livraison~; surveillance des \omterm{def:rc:operational-risk-and-rogue-trading:sod}{congés obligatoires consécutifs} (block leave).
\end{build}

\begin{omsources}
\item House of Commons, \emph{Hansard}, 18 juillet 1995, déclaration sur le rapport du Board of Banking Supervision
sur Barings.
\item Financial Services Authority, Final Notice to UBS AG, 25 novembre 2012.
\item Commodity Futures Trading Commission, release 4144-98, 11 mai 1998.
\item Board of Governors of the Federal Reserve System, ordonnance contre Daiwa Bank, 2 octobre 1995.
\item Allied Irish Banks, rapport annuel sur le Form 20-F pour 2001 (Allfirst).
\item Société Générale, General Inspection Department, \emph{Mission Green}, rapport du 20 mai 2008.
\item Basel Committee on Banking Supervision, \emph{Basel~III: Finalising post-crisis reforms}, décembre
2017 (risque opérationnel).
\end{omsources}

\section{Exercices}

\begin{exercise}[$\star$]\label{exo:rc:operational-risk-and-rogue-trading:1}
Calculer le BIC d'une banque dont l'indicateur d'activité s'élève à 10 milliards d'euros.
\end{exercise}

\begin{exercise}[$\star$]\label{exo:rc:operational-risk-and-rogue-trading:2}
Pourquoi l'ILM est-il égal à un lorsque la composante de perte égale le BIC~?
\end{exercise}

\begin{exercise}[$\star$]\label{exo:rc:operational-risk-and-rogue-trading:3}
Quelle règle aurait signalé les transactions sur ETF à règlement différé du système UBS~?
\end{exercise}

\begin{exercise}[$\star\star$]\label{exo:rc:operational-risk-and-rogue-trading:4}
Pourquoi prêter 300 millions de livres sterling contre 31 millions de livres sterling de collatéral constitue-t-il un signal de contrôle~?
\end{exercise}

\begin{exercise}[$\star\star$]\label{exo:rc:operational-risk-and-rogue-trading:5}
Pourquoi exiger deux règles réduit-il les fausses alertes bien davantage que les détections~?
\end{exercise}

\begin{exercise}[$\star\star$]\label{exo:rc:operational-risk-and-rogue-trading:6}
Comment les \omterm{def:rc:operational-risk-and-rogue-trading:sod}{congés obligatoires consécutifs} (block leave) exposent-ils une position dissimulée~?
\end{exercise}

\begin{exercise}[$\star\star\star$]\label{exo:rc:operational-risk-and-rogue-trading:7}
\emph{Programmation.} Vérifier que le journal d'audit détecte la modification du prix d'une transaction, et expliquer pourquoi
la réécriture de toute la chaîne après une modification serait tout de même détectée si le dernier hachage est conservé ailleurs.
\end{exercise}

\begin{exercise}[$\star\star\star$]\label{exo:rc:operational-risk-and-rogue-trading:8}
\emph{Trouver la faille.} «~Le P\&L de notre trader est régulier et élevé avec une VaR faible~; c'est notre meilleur élément, et
nous devrions relever ses limites.~»
\end{exercise}

\section{Problème~: janvier 2008}

\begin{problem}[{Problème du week-end --- dimensionner une position dissimulée}]\label{pb:rc:operational-risk-and-rogue-trading:1}
Le rapport d'inspection de Société Générale a établi qu'un trader a constitué une position acheteuse de 49~milliards d'euros
sur des contrats à terme sur indices, entre le 2 et le 18 janvier 2008, découverte le 20~; le dénouement de cette position, du 21 au 23 janvier,
a coûté 6.4~milliards d'euros, soit 4.9~milliards d'euros nets des 1.5~milliard gagnés par ses positions en 2007.

\medskip\noindent\textbf{Partie I --- L'arithmétique.}
\begin{enumerate}
\item Donner le mouvement défavorable moyen qu'impliquent la perte et la position.
\item Quelle position la même perte impliquerait-elle pour un mouvement défavorable moyen de 5~\%~?
\item Pourquoi le dénouement lui-même a-t-il fait évoluer le marché contre la banque~?
\item En supposant une volatilité annuelle de 25~\% pour les actions, quelle VaR à un jour au niveau de confiance de 99~\%
la position aurait-elle affichée~?
\item Pourquoi la VaR de la banque n'a-t-elle pas fait apparaître la position~?
\end{enumerate}

\medskip\noindent\textbf{Partie II --- Les contrôles.}
\begin{enumerate}[resume]
\item Quel type d'\omterm{def:rc:operational-risk-and-rogue-trading:rogue}{opérations fictives} dissimule une large position directionnelle~?
\item Laquelle des règles du chapitre permettrait de les repérer~?
\item Qu'auraient révélé les confirmations auprès des contreparties~?
\item Quel signal de trésorerie ou d'appel de marge aurait dû apparaître~?
\item Quel rôle jouent les \omterm{def:rc:operational-risk-and-rogue-trading:sod}{congés obligatoires consécutifs} (block leave)~?
\end{enumerate}

\medskip\noindent\textbf{Partie III --- La détection en pratique.}
\begin{enumerate}[resume]
\item Donner les taux de détection et de fausses alertes des règles combinées sur le registre des opérations du chapitre.
\item Combien d'alertes par an un desk de cette taille produirait-il, et qui les instruit~?
\item Pourquoi des alertes sont-elles clôturées sans avoir été résolues~?
\item Comment le journal d'audit aide-t-il une enquête a posteriori~?
\item Quels éléments faudrait-il ajouter aux règles pour un desk de futures~?
\end{enumerate}

\medskip\noindent\textbf{Partie IV --- Jugement.}
\begin{enumerate}[resume]
\item Pourquoi les traders fraudeurs commencent-ils en général par dissimuler des pertes plutôt que par rechercher des gains~?
\item Quel contrôle constitue la première ligne de défense qui l'aurait repéré~?
\item Comment les fonds propres au titre du \omterm{def:rc:operational-risk-and-rogue-trading:oprisk}{risque opérationnel} doivent-ils refléter de tels événements~?
\item Énoncer le \emph{résultat nommé}~: le mouvement moyen qu'impliquent la perte déclarée et la position, et le
premier contrôle de la chaîne qui l'aurait repéré.
\item En une phrase~: qu'est-ce qui rend un contrôle efficace~?
\end{enumerate}
\end{problem}

\section{Questions d'entretien}

\begin{interviewq}[$\star$ \iqroles{risque, banque}]\label{iq:rc:operational-risk-and-rogue-trading:1}
Définissez le \omterm{def:rc:operational-risk-and-rogue-trading:oprisk}{risque opérationnel}. En quoi se distingue-t-il du risque de marché et du risque de crédit~?
\end{interviewq}

\begin{interviewq}[$\star\star$ \iqroles{risque}]\label{iq:rc:operational-risk-and-rogue-trading:2}
Quels contrôles préviennent la \omterm{def:rc:operational-risk-and-rogue-trading:rogue}{fraude de trader} (rogue trading)~? Lesquels ont failli chez Barings~?
\end{interviewq}

\begin{interviewq}[$\star\star$ \iqroles{développeur}]\label{iq:rc:operational-risk-and-rogue-trading:3}
Concevez un système de surveillance des \omterm{def:rc:operational-risk-and-rogue-trading:rogue}{opérations fictives}.
\end{interviewq}

\begin{interviewq}[$\star\star$ \iqroles{trader}]\label{iq:rc:operational-risk-and-rogue-trading:4}
Le P\&L d'un collègue est régulier, important et inexpliqué par ses positions. Que faites-vous~?
\end{interviewq}

\begin{interviewq}[$\star\star\star$ \iqroles{banque, risque}]\label{iq:rc:operational-risk-and-rogue-trading:5}
Comment les fonds propres au titre du \omterm{def:rc:operational-risk-and-rogue-trading:oprisk}{risque opérationnel} sont-ils calculés selon l'approche standard de Bâle, et quelles en sont les faiblesses~?
\end{interviewq}

\begin{interviewq}[$\star\star\star$ \iqroles{risque}]\label{iq:rc:operational-risk-and-rogue-trading:6}
Pourquoi les alertes de surveillance ne parviennent-elles pas à empêcher les opérations non autorisées, et comment y remédieriez-vous~?
\end{interviewq}
